\section{Empirical Strategy}
\label{sec:method}

\subsection{The estimand}

For each quantile $\tau$ and horizon $h$ we estimate, by local projection
\citep{jorda2005}, the response of the conditional return quantile to the
liquidation shock:
\begin{equation}
\label{eq:qlp}
\QR_{\tau}\!\left(r_{t+1 \to t+1+h} \mid \mathcal{F}_t\right)
 = \beta_h(\tau)\, L_{t-1}
 + \delta_h(\tau)\, \left(L_{t-1} \times \texttt{oi\_high}_t\right)
 + \gamma_h(\tau)' X_t ,
\end{equation}
where $r_{t+1 \to t+1+h}$ is the cumulative ETH return over the $h+1$ hours from $t+1$ to $t+1+h$ (so $h=0$ is the one-hour-ahead return),
$L_{t-1}$ the lagged log-liquidation shock, and $X_t$ the controls of
Section~\ref{sec:method-controls}. The object of interest is
$\beta_h(\tau)$ across $\tau \in \{0.01, 0.05, 0.10, 0.50, 0.90, 0.95\}$
and $h \in \{0, \dots, 24\}$ hours (a nine-quantile grid is reported in the
appendix). Estimation is by quantile regression \citep{koenker1978}; the
estimator and its bootstrap inference follow \citet{han2024}. The reported
grid stops at $\tau = 0.95$; the tail-inference and mirror set additionally
estimates $\tau = 0.99$, so that the one-percent upper tail has a symmetric
counterpart to the one-percent lower tail. Both grids are fixed in the locked
configuration file, so $0.99$ is estimated but not carried into the reported
panels.

The quantile approach is necessary rather than cosmetic, for two reasons.
First, the hypothesis under test, leverage-cycle amplification
\citep{geanakoplos2010,brunnermeier2009}, is a statement about the
\emph{tails} of the return distribution, not its mean. Second, the mean is
empirically silent here: an OLS local projection of returns on the shock
finds nothing at any horizon (point estimate $\approx \OlsLpBetaHzero$ on
impact, HAC $p = \OlsLpPHzero$). This is the Growth-at-Risk logic of
\citet{adrian2019}
transposed to high-frequency crypto returns: tails can respond where the
mean is blind. We note at the outset a tension that
Section~\ref{sec:deconstruction} resolves. \citet{adrian2019} interpret
asymmetric tail sensitivity through conditional skewness;
\citet{carriero2024} show that symmetric conditional distributions, under joint shifts in mean and variance, can produce the same apparent pattern; and \citet{brownlees2021} find that a
plain GARCH rivals quantile-regression Growth-at-Risk out of sample.
Distinguishing those readings is the heart of this paper.

\subsection{Shock definition and timing}
\label{sec:method-shock}

The canonical shock is \emph{raw}: $L_{t-1} = \ln(1 +
\text{liq}^{\text{USD}}_{t-1})$, one-hour lagged
(Section~\ref{sec:data}). Timing is read strictly. ``$h=0$'' is a
one-hour-ahead prediction; within the bucket at $h=-1$, prices and
liquidations co-move mechanically (the reverse-causality zone of Fact 2);
only $h \le -2$ would be a true lead. All claims are therefore
predictive associations.

A second, \emph{orthogonalized} shock (the residual of an OLS regression
of the log flow on contemporaneous and lagged BTC returns, then lagged)
is reserved for the cross-asset placebo (Test A of the robustness battery),
in the ex-ante residualization spirit of \citet{borusyak2022}. The two
definitions preserve sign and significance at every horizon and agree to
within a few percent from six hours out, but diverge at short horizons (the
impact tail coefficient is about 60\% larger under orthogonalization). The
choice between raw and orthogonalized is thus an identification convention
that leaves the symmetry finding and the bound intact, not a claim of
numerical coincidence at every horizon; residualizing at the quantile (rather than by OLS)
is an acknowledged refinement left for future work.

\subsection{Controls and the identification logic}
\label{sec:method-controls}

The regressor set is locked:
$\{L_{t-1},\; L_{t-1}\times\texttt{oi\_high}_t,\; \texttt{oi\_high}_t,\;
r^{\text{BTC}}_t,\; \sigma^{7d}_t,\; \text{funding}_t,\; \text{basis}_t\}$.
The realized-volatility control $\sigma^{7d}_t$ (168-hour rolling standard
deviation) is not a nuisance term: it is central to the identification.
Liquidations raise volatility, and volatility widens \emph{both}
tails symmetrically; the leverage-cycle prediction is specifically
\emph{downside} amplification. Holding $\sigma^{7d}_t$ fixed conditions
$\beta_h(\tau)$ on the \emph{predetermined} volatility state, removing the
confound between the shock and the pre-existing volatility regime. It does
not remove the shock's own forward-volatility response
(Section~\ref{sec:volchannel}). But that response is symmetric, so it cannot
manufacture asymmetry; it is the symmetric residual that the placebo and
the paired tail tests identify directly. The asymmetry tests, not this
control alone, are therefore what separate a real downside shape from a
symmetric scale effect. That a symmetric conditional distribution can
masquerade as downside risk is the demonstration of \citet{carriero2024},
there through joint shifts in mean and variance; the scale-versus-shape
distinction is what quantile regression makes testable
\citep{koenker1982,machado2000}. Realized volatility is preferred to a
GARCH or implied measure because it is observable, model-free, and
predetermined; model-based volatility enters only inside the placebo DGPs
of Section~\ref{sec:deconstruction}, where its model-dependence is the
point.

\subsection{Inference}
\label{sec:method-inference}

Inference proceeds on two tiers. At central quantiles, kernel standard
errors (Hall--Sheather bandwidth) are retained. At tail quantiles
($\tau \in \{0.01, 0.05, 0.95, 0.99\}$), the primary inferential object is
the percentile confidence interval from a moving-block bootstrap with
24-hour blocks \citep{fitzenberger1998,politis1994,gregory2018}, for two
reasons. First, the cumulative left-hand side is overlapping by
construction, so residuals are serially correlated at every $h > 0$
\citep{montiel2021,lusompa2023}; blocks of one day, the maximal horizon,
preserve that dependence. Second, at $\tau = 0.01$ roughly 406
observations sit below the conditional quantile, and the extremal quantile
regression literature shows the asymptotic distribution is non-Gaussian
that far in the tail \citep{chernozhukov2005,chernozhukov2011}, making
kernel inference unreliable: on the exact headline specification, the
bootstrap-to-kernel standard-error ratio at $\tau = 0.01$ ranges from
\SeRatioMin{} to \SeRatioMax{} across horizons
(Appendix Table~\ref{tab:seratio}). All bootstrap results use $1{,}000$
replications with a deterministic multi-level seeding scheme. The MBB block
length does not drive the results: re-running the tail inference at 12-, 36-, and
48-hour blocks leaves every interval and every minimum detectable effect
essentially unchanged (Appendix Table~\ref{tab:blocksens}).

\subsection{The diagnostic battery}
\label{sec:method-battery}

The core methodological contribution is a battery of diagnostics. Each one is
designed to isolate a single confound of the naive reading. What survives all
of them is what we are willing to call a real signal.

\begin{table}[t]
\centering
\caption{Diagnostic tests, one per confound.}
\label{tab:battery}
{\small
\begin{tabular}{p{0.27\textwidth}p{0.21\textwidth}p{0.42\textwidth}}
\toprule
Diagnostic & Confound isolated & Logic \\
\midrule
Symmetric sign-flip placebo & volatility (scale) & does a zero-skew world
with the real volatility path reproduce $\beta(0.01)$ and the left--right
gap? \citep{wu1986,liu1988,davidson2008,carriero2024} \\
Dual permutation nulls & overlap artifact & destroy the shock--return link
(circular shift; innovation reshuffle), keep the cumulative LHS: any
surviving $\beta$ is mechanical \citep{hodrick1992,boudoukh2008,valkanov2003} \\
Tail-exceedance regressions & (the predictor + the symmetry test) & does
the shock predict future tail violations, and equally on both sides?
\citep{engle2004,balcilar2017} \\
Standardized skew tests & scale noise & standardize by predetermined volatility, then
test down-minus-up and the third moment \citep{neuberger2021,farago2023} \\
BTC-outcome placebo & asset specificity & re-run the full QLP with BTC as
the outcome: does the signature replicate on an asset without the
ETH-collateral channel? \\
Equivalence bounds (MDE/TOST) & power & convert ``not found'' into ``we
exclude effects larger than $X$''
\citep{altman1995,ioannidis2017,lakens2017,riesthuis2024} \\
Out-of-sample pinball test & in-sample overfit & does liquidation
information improve rolling-origin tail forecasts against nested and
GARCH benchmarks? \citep{diebold1995} \\
Planted-cascade positive control & power (false negatives) & does the
battery recover a sign-specific downside asymmetry planted at known scale,
and stay silent below it? \\
\bottomrule
\end{tabular}}
\end{table}

Table~\ref{tab:battery} summarizes the battery; applications and results
are in Sections~\ref{sec:deconstruction}--\ref{sec:survives}. The
equivalence step requires a pre-specified smallest effect size of interest
(SESOI). We set it at a one-percentage-point difference in tail-violation
probability between the downside and upside over the span of the shock,
translated into coefficient units under three calibrations of that span
(interquartile, median-to-95th, and 10th-to-90th percentile of the nonzero
shock); the strictest yields \SesoiStrict{} per unit of the log shock. The
one-point threshold is decisional, not statistical. It is the smallest
down-versus-up tail miscalibration a Value-at-Risk desk would act on, not the
smallest the sample can resolve \citep{lakens2017}. One percentage point is
the order at which a shift in the realized 1\% exceedance rate moves a model
across the Basel traffic-light boundary over the standard 250-day backtesting
window. That window is the institutional unit at which tail calibration is
judged by mandate, and we anchor the threshold to it rather than to the far
finer statistical resolution of a forty-thousand-hour sample.

\subsection{Multiple testing and the locked specification}

Bonferroni corrections guard the grid searches: a $\times 12$ family over
the reported $\tau \times h$ cells and a $\times 6$ family over the
quantile grid. The correction family is the set of reported cells itself, the
$\tau \times h$ grid and the quantile grid enumerated in the locked
configuration file before the diagnostic phase, so the family is fixed by the
specification rather than chosen after seeing the diagnostics. The direct shock effect survives both families in full
(12/12 and 6/6); the leverage interaction $\delta_h(\tau)$ fails all six
(0/6). We carry that null forward as it stands, since the interaction was the
specification's embodiment of the leverage-cycle prediction. We read it as a
\emph{detectability} null, not a tight bound (the interaction is imprecisely
estimated; Section~\ref{sec:apparent-crack}), and the paper's quantitative
bounds rest on the asymmetry objects, not on this count. The
corrections guard against noise, not against artifacts, which is exactly why
the battery of Table~\ref{tab:battery} exists.

The specification was locked before the diagnostic phase: every parameter
(quantile grids, horizons, controls, block length, replication counts,
seeds) is fixed in a single configuration file, and every diagnostic is a
flagged auxiliary analysis, never a modification of
Equation~\eqref{eq:qlp}. The full pipeline, from panel construction through
estimation, robustness, diagnostics, and every number, table, and figure in
this paper, regenerates from public data with a one-command replication
package; local smoke runs and the canonical compute environment reproduce
point estimates to three decimals.